% !TeX root = ../main.tex
% Add the above to each chapter to make compiling the PDF easier in some editors.
\chapter{Background and Fundamentals}
\label{chap:background}

This chapter introduces the fundamental concepts required for
understanding the research presented in this dissertation.
First, the foundations of real-time systems and scheduling
theory are discussed, including timing constraints, scheduling
algorithms, and real-time operating systems. Afterwards, the
concept of mixed-criticality systems is introduced together
with the challenges associated with resource sharing and
adaptability in modern embedded platforms.
Subsequently, distributed real-time systems and modern
automotive computing architectures are analyzed, including
heterogeneous systems, edge computing, and task offloading
concepts. Finally, the chapter presents the networking
technologies and communication mechanisms commonly used in
distributed cyber-physical systems with real-time requirements.

% ============================================================
\section{Real-Time Systems}
\label{sec:rtsystems}

Real-time systems are computing systems in which correctness
depends not only on the logical correctness of the produced
results, but also on the time at which these results are
generated~\cite{buttazzo2011}. In such systems, timing
constraints are treated as first-class requirements and must
be satisfied in order to guarantee correct system behavior.
Real-time systems are commonly classified into hard real-time,
firm real-time, and soft real-time systems depending on the
consequences of missing timing deadlines. In hard real-time
systems, missing a deadline may lead to catastrophic failures,
making deterministic timing guarantees essential. Firm
real-time systems occupy an intermediate position: a missed
deadline renders the result valueless, but unlike hard
real-time systems, individual misses do not necessarily
trigger a system failure. Soft real-time systems, on the other
hand, can tolerate occasional deadline misses with graceful
degradation of quality of service.

Each task in a real-time system is typically characterized by
a set of parameters: its release time (or arrival time), an
execution time, a relative deadline, and---for periodic
tasks---a period. The Worst-Case Execution Time~(WCET) is a
particularly critical parameter, as it represents an upper
bound on the time a task may require under any input and
execution condition. Deriving tight and safe WCET estimates
is a research challenge in its own right, with techniques
ranging from static analysis~\cite{wilhelm2008} to
measurement-based approaches~\cite{engblom2002}. The
inherent trade-off between safety (overestimation) and
efficiency (tight bounds) is especially pronounced in
heterogeneous and many-core platforms, where cache effects,
pipeline speculation, and memory contention increase
analysis complexity.

Modern cyber-physical systems such as automotive platforms,
industrial controllers, robotics systems, and autonomous
vehicles increasingly depend on strict timing guarantees while
simultaneously requiring higher computational flexibility and
adaptability~\cite{baunach2018,sdv2024}. This combination
introduces significant challenges for system design and
resource management.

% ------------------------------------------------------------
\subsection{Scheduling Theory}
\label{subsec:scheduling}

Scheduling is one of the core concepts in real-time systems
and refers to the process of determining the execution order
of tasks while satisfying timing constraints and resource
limitations. Classical real-time scheduling approaches include
fixed-priority scheduling algorithms such as Rate Monotonic
(RM), as well as dynamic-priority approaches such as
Earliest-Deadline-First (EDF).

Liu and Layland~\cite{liu1973} introduced one of the
foundational scheduling models for periodic real-time tasks,
establishing schedulability bounds for RM and EDF scheduling.
EDF scheduling is particularly relevant due to its optimality
properties for uniprocessor systems under specific assumptions:
it is known to be optimal in the sense that if any algorithm
can schedule a task set on a single processor, EDF can also
schedule it. Fixed-priority schemes such as RM, while not
theoretically optimal, offer simpler analysis and more
predictable behavior in the presence of task interactions,
making them dominant in practice.

Schedulability analysis is the formal process of determining
whether a task set will always meet its deadlines under a
given scheduling policy. Response Time Analysis~(RTA) for
fixed-priority systems~\cite{joseph1986} and utilization-based
bounds for EDF~\cite{liu1973} are well-established analytical
tools. However, extending such analyses to account for
blocking from shared resources, interrupt handling, and
inter-task communication requires additional models such as
the Sporadic Server or the Stack Resource Policy~\cite{baker1991}.

As modern embedded systems evolved towards multicore and
heterogeneous architectures, scheduling approaches also became
more complex. Multiprocessor scheduling introduces additional
challenges including task migration overhead, synchronization,
resource sharing, cache interference, and communication
latency~\cite{megel2010}. Broadly, multiprocessor scheduling
strategies are classified into partitioned, global, and
semi-partitioned (or clustered) approaches. Partitioned
scheduling assigns each task permanently to one processor and
applies uniprocessor analysis independently, at the cost of
possible inefficiency due to bin-packing constraints. Global
scheduling allows tasks to migrate freely across processors,
improving utilization but introducing runtime migration
overhead and making formal analysis significantly harder.
Semi-partitioned and cluster-based approaches attempt to
balance these trade-offs~\cite{davis2011}.

% ------------------------------------------------------------
\subsection{Real-Time Operating Systems}
\label{subsec:rtos}

Real-Time Operating Systems (RTOSs) provide the software
infrastructure required for executing time-critical tasks
while maintaining deterministic system behavior. Compared to
general-purpose operating systems, RTOSs prioritize
predictability, bounded latency, and deterministic scheduling
behavior over throughput optimization.

Common RTOS functionalities include task scheduling,
inter-process communication, synchronization primitives,
interrupt handling, and memory management. Depending on the
system requirements, RTOSs may support static scheduling,
dynamic scheduling, partitioning, or mixed-criticality
execution models. Well-known examples include
OSEK/VDX~\cite{osek2005} and AUTOSAR Classic Platform, which
are widely deployed in automotive embedded systems, as well as
more general-purpose RTOSs such as FreeRTOS and VxWorks.

The AUTOSAR Adaptive Platform~\cite{autosar2023} represents a
more recent development targeted at high-performance
automotive computing nodes. Unlike the Classic Platform, which
is oriented towards microcontrollers with static task
configurations, the Adaptive Platform supports dynamic
application deployment, service-oriented communication, and
POSIX-compliant execution environments---all while maintaining
defined timing and isolation guarantees. This architectural
shift is particularly relevant for the integration of
mixed-criticality workloads and dynamic task management in
modern vehicle platforms.

In distributed and heterogeneous systems, RTOSs additionally
play an important role in enabling communication, coordination,
and resource sharing between different hardware nodes. The
ability to support runtime adaptation and dynamic task
migration increasingly represents an important requirement for
modern distributed embedded systems. Hypervisor-based
virtualization is increasingly used to consolidate multiple
RTOS instances or mixed-criticality partitions on a single
hardware platform~\cite{heiser2020}, providing spatial and
temporal isolation with lower hardware costs.

% ============================================================
\section{Mixed-Criticality Systems}
\label{sec:mcs}

Mixed-criticality systems~(MCS) are systems in which
applications with different levels of criticality coexist on
shared hardware resources~\cite{vestal2007,burns2017}. Such
systems are especially relevant in domains including
automotive systems, avionics, robotics, and industrial
automation, where safety-critical and non-safety-critical
applications execute simultaneously within the same platform.

Traditionally, critical systems relied on strong spatial and
temporal partitioning mechanisms in order to isolate workloads
with different certification requirements. While this approach
improves safety and predictability, it may also lead to low
resource utilization and reduced flexibility.

Vestal~\cite{vestal2007} introduced one of the most influential
mixed-criticality models, where tasks are associated with
different Worst-Case Execution Time (WCET) estimations
depending on their assurance level. Specifically, a
high-criticality task carries a pessimistic (safe) WCET
estimate for use under high-assurance conditions, and a
more optimistic estimate for use when the system operates
in a nominal low-load mode. Since then, multiple scheduling
approaches have been proposed to support dynamic adaptation
and graceful degradation under overload conditions. The
mixed-criticality model of Baruah et al.~\cite{baruah2011}
extended Vestal's original formulation to a formal
schedulability framework under EDF, providing a foundation
for a significant body of subsequent research.

% ------------------------------------------------------------
\subsection{Criticality Levels and Standards}
\label{subsec:criticality_standards}

In safety-critical domains, criticality levels are often
defined by domain-specific standards. In avionics, the
DO-178C standard defines Design Assurance Levels (DALs)
A through E for software, where DAL~A represents the
highest assurance requirement for functions whose failure
could cause catastrophic consequences~\cite{do178c2011}. The
corresponding hardware standard DO-254 applies analogous
levels to electronic hardware design. In automotive systems,
ISO~26262 defines Automotive Safety Integrity Levels
(ASILs) ranging from ASIL~A to ASIL~D, with ASIL~D
representing the most stringent requirements for functions
whose failure could lead to severe injuries or
fatalities~\cite{iso26262_2018}. For general functional
safety across industrial domains, IEC~61508 defines Safety
Integrity Levels (SILs) 1 through 4.

These standards impose requirements not only on the software
algorithms themselves but also on the processes, tools,
verification methods, and evidence used during development.
When components of different criticality share hardware
resources, demonstrating compliance becomes significantly
more complex, as the behavior of lower-criticality components
may potentially interfere with higher-criticality ones.
Achieving compositionality---the ability to verify system
properties from verified component properties---is a central
challenge in the design of mixed-criticality platforms.

% ------------------------------------------------------------
\subsection{Scheduling and Resource Management for MCS}
\label{subsec:mcs_scheduling}

Modern mixed-criticality systems increasingly explore adaptive
resource management techniques in order to improve efficiency
while preserving timing guarantees~\cite{adaptive2023}.
In the standard mixed-criticality scheduling model, the
system can transition between a low-criticality (LO) mode
and a high-criticality (HI) mode. When the observed execution
time of a high-criticality task exceeds its LO-mode WCET
estimate, the system switches to HI mode, abandoning or
reducing the service of low-criticality tasks in order to
guarantee that all high-criticality deadlines are met.

Several extensions to this basic model have been proposed.
The AMC (Adaptive Mixed Criticality) scheduler~\cite{burns2013}
uses per-task mode switching rather than a global transition,
allowing finer-grained adaptation. Elastic scheduling
approaches attempt to dynamically adjust task parameters
(periods, execution budgets) in response to load changes
while maintaining feasibility constraints. Research on
\emph{service degradation} policies addresses the question of
how best to reduce the service of lower-criticality tasks
during overload, recognizing that an abrupt shutdown of
lower-criticality functions may itself introduce safety risks.

Recent research has also investigated stronger isolation
mechanisms and heterogeneous execution models for
mixed-criticality platforms~\cite{hoornaert2024,tinto2024}.
Cache coloring and cache partitioning have emerged as
important mechanisms for reducing cache interference between
tasks of different criticality levels executing on shared-cache
multicore platforms~\cite{mancuso2013}. Memory bandwidth
regulation and the use of memory protection units~(MPUs) or
memory management units~(MMUs) further contribute to
spatial isolation guarantees.

In distributed cyber-physical systems such as modern
automotive platforms, mixed-criticality introduces additional
complexity due to communication delays, distributed resource
sharing, synchronization constraints, and heterogeneous
hardware architectures. These systems must simultaneously
maintain timing guarantees, fault tolerance, and runtime
adaptability across multiple interconnected nodes. The
integration of mixed-criticality scheduling with distributed
resource management and task migration constitutes a central
open challenge addressed in this dissertation.

% ============================================================
\section{Distributed Real-Time Systems}
\label{sec:distributed}

Distributed real-time systems consist of multiple computing
nodes that cooperate in order to execute time-critical tasks
while communicating over a network infrastructure. Such
systems are commonly used in domains including automotive
electronics, industrial automation, aerospace systems, and
large-scale cyber-physical systems. In contrast to
uniprocessor or shared-memory multiprocessor systems,
distributed systems must cope with the inherent uncertainty
of network communication, including variable transmission
delays, packet losses, and clock synchronization errors,
all while satisfying end-to-end timing constraints.

Modern automotive systems increasingly evolve from
domain-oriented architectures composed of numerous Electronic
Control Units~(ECUs) towards centralized and zonal
architectures based on high-performance computing
platforms~\cite{sdv2024,automotivehpc2023}. This transition
is motivated by the growing computational requirements of
software-defined functionalities, autonomous driving systems,
and advanced driver assistance systems~(ADAS). A zonal
architecture groups ECUs by physical location in the vehicle
rather than by function, with high-bandwidth Ethernet backbones
connecting zone controllers to central compute nodes. This
consolidation reduces wiring complexity and weight while
enabling more powerful centralized processing, but also
concentrates risk and demands stronger mechanisms for
isolation and fault containment.

Distributed architectures provide several advantages,
including scalability, fault tolerance, workload
distribution, and resource sharing. However, they also
introduce challenges associated with synchronization,
communication latency, scheduling coordination, and
distributed resource management.

% ------------------------------------------------------------
\subsection{End-to-End Timing Analysis}
\label{subsec:e2e}

A fundamental challenge in distributed real-time systems is
the analysis of end-to-end latency for processing chains
that span multiple nodes. A \emph{cause-effect chain} (or
task chain) is a sequence of tasks executing on potentially
different nodes, where each task consumes the output of the
previous one. The end-to-end latency of such a chain is
determined not only by the execution times of individual
tasks, but also by their activation patterns, scheduling
decisions on each node, and the communication delays
incurred at each link.

The problem of bounding end-to-end latency has been studied
through several analytical frameworks. The work of
Dürr et al.~\cite{durr2019} provides a systematic
characterization of end-to-end latency semantics,
distinguishing between \emph{reaction time} (the worst-case
delay from an external event to the completion of the last
task in the chain) and \emph{data age} (the maximum
staleness of data observed by the final consumer). These
two metrics capture different aspects of real-time
correctness relevant to different application domains;
for example, control systems are typically more sensitive to
data age, whereas event-triggered safety responses demand
short reaction times.

% ------------------------------------------------------------
\subsection{Heterogeneous and Edge Computing Platforms}
\label{subsec:hetero_edge}

Modern distributed embedded systems frequently combine
heterogeneous computing resources including CPUs, GPUs,
FPGAs, and specialized accelerators. Such heterogeneous
platforms can improve performance and energy efficiency but
also increase system complexity and scheduling challenges.
The diversity of instruction set architectures, memory
hierarchies, and programming models across these processing
elements means that tasks cannot in general be freely
migrated without adaptation of the application binary or
execution environment.

The emergence of edge and Multi-access Edge Computing~(MEC)
architectures additionally enables the distribution of tasks
between embedded devices, edge servers, and cloud
infrastructure~\cite{mec2017}. This paradigm supports
workload offloading, adaptive deployment, and resource
elasticity while reducing latency compared to centralized
cloud computing approaches. In latency-sensitive
cyber-physical applications, edge computing plays a
particularly important role by enabling local processing
of sensor data with sub-millisecond response times that
would be infeasible with round-trips to a remote cloud.

Task offloading and workload distribution mechanisms are
closely related to orchestration concepts commonly used in
cloud and Internet-of-Things (IoT) environments. In such
systems, orchestration mechanisms coordinate deployment,
resource allocation, monitoring, and execution management
across multiple distributed nodes. Container-based
virtualization and lightweight orchestration frameworks,
originally developed for cloud-native applications, are
increasingly being adapted for edge and embedded
contexts~\cite{morabito2018}, though the real-time and
safety guarantees required in cyber-physical systems
impose constraints that are not natively addressed by
mainstream cloud orchestration tools.

% ------------------------------------------------------------
\subsection{Task Migration and Runtime Adaptation}
\label{subsec:taskmig}

Task migration refers to the process of transferring a task
or workload from one execution node to another during system
operation. Migration mechanisms may be used for fault
tolerance, load balancing, resource optimization, energy
management, or adaptive scheduling purposes.

Existing migration approaches range from static failover
mechanisms to fully dynamic runtime migration strategies.
However, implementing migration in real-time systems
introduces several challenges, including state transfer,
synchronization overhead, migration latency, and timing
predictability~\cite{chen2012,megel2010}. The migration
process itself can be decomposed into several phases:
(i)~the \emph{pre-migration} phase, in which the source
node prepares and serializes the task state; (ii)~the
\emph{transfer} phase, during which state data is
transmitted to the destination node over the interconnect;
and (iii)~the \emph{post-migration} phase, in which the
destination node initializes execution from the transferred
state. Each of these phases introduces latency that must
be accounted for in the schedulability analysis of the
migrating task as well as of any tasks that share resources
with it.

Checkpoint-based migration approaches periodically save
a task's execution state to stable storage or shared memory,
enabling resumption from the most recent checkpoint after
a failure or migration event~\cite{elnozahy2002}. While
this reduces the cost of the transfer phase by transmitting
incremental state differences, it introduces regular runtime
overhead for checkpointing and may not be compatible with
hard real-time timing guarantees unless the checkpointing
overhead is explicitly bounded and included in schedulability
analysis.

In heterogeneous distributed systems, migration becomes even
more complex due to architectural differences between hardware
platforms, communication constraints, and execution
environment compatibility. Binary incompatibility between
different instruction set architectures requires either
maintaining multiple compiled variants of each migratable
task, recompilation at migration time, or the use of
intermediate representations or virtual machines that
abstract the underlying hardware. The overhead introduced
by such compatibility layers must be carefully evaluated
against the latency and correctness requirements of the
target application.

% ============================================================
\section{Networks and Communication in Real-Time Systems}
\label{sec:networks}

Communication networks represent a critical component of
distributed cyber-physical systems. In real-time systems,
communication mechanisms must provide deterministic behavior,
bounded latency, synchronization support, and fault tolerance.

Traditional automotive systems relied heavily on fieldbus
technologies such as Controller Area Network~(CAN),
FlexRay, and Local Interconnect Network~(LIN). CAN provides
a broadcast, message-priority-based medium access scheme
with well-understood worst-case latency bounds, but its
bandwidth is limited to 1~Mbit/s (or up to 8~Mbit/s with
CAN~FD). FlexRay offers higher bandwidth and a time-triggered
communication model that provides stronger temporal
determinism, but at higher cost and complexity. LIN is used
for low-bandwidth sensor and actuator networks where cost
is paramount. These technologies collectively supported the
domain-oriented ECU architectures of previous vehicle
generations.

% ------------------------------------------------------------
\subsection{Time-Sensitive Networking}
\label{subsec:tsn}

More recent architectures increasingly incorporate
Ethernet-based technologies and Time-Sensitive Networking~(TSN)
mechanisms in order to support higher bandwidth requirements
and improved scalability~\cite{ieee8021q2022}. TSN extends
standard IEEE~802.3 Ethernet with a suite of standards
developed by the IEEE~802.1 working group, collectively
providing mechanisms for traffic prioritization, time
synchronization, bounded latency, and deterministic
communication.

Key TSN standards include IEEE~802.1AS for generalized
precision time protocol~(gPTP)-based clock
synchronization~\cite{ieee8021as2020}, IEEE~802.1Qbv for
time-aware shaping~(TAS), and IEEE~802.1Qbu/802.3br for
frame preemption, which allows high-priority frames to
interrupt the transmission of lower-priority frames. The
Credit-Based Shaper~(CBS) defined in IEEE~802.1Qav provides
bandwidth reservation for audio-video bridging~(AVB) streams.
Together, these mechanisms allow TSN networks to support
multiple traffic classes with heterogeneous latency and
bandwidth requirements on a shared infrastructure.

TSN capabilities are especially relevant for distributed
real-time systems where tasks executing on different nodes
require predictable communication behavior. In the context
of task migration, TSN networks can provide bounded-latency
channels for state transfer, coordination messages, and
monitoring traffic, enabling more predictable migration
latency bounds than best-effort networks. However, the
interaction between TSN traffic scheduling, task scheduling
on computing nodes, and the timing constraints of
distributed applications remains an active area of
research~\cite{craciunas2016}.

% ------------------------------------------------------------
\subsection{Service-Oriented Communication and Middleware}
\label{subsec:middleware}

In addition to transport-layer networking technologies, the
communication middleware used in distributed real-time
systems plays a significant role in system integration and
runtime behavior. The AUTOSAR Adaptive Platform adopts a
service-oriented communication model based on the SOME/IP
protocol~\cite{autosar2023}, which supports dynamic service
discovery and method-call semantics over standard IP
networks. The Data Distribution Service~(DDS) middleware
standard~\cite{omgdds2015} provides a publish-subscribe
communication model with Quality-of-Service~(QoS) policies
governing reliability, latency budgets, deadlines, and
resource limits, and has seen growing adoption in
autonomous and robotic systems through its use as the
underlying transport in the Robot Operating System~2
(ROS2)~\cite{maruyama2016}.

In distributed task migration and orchestration scenarios,
network infrastructure additionally plays an important role
in state transfer, coordination, synchronization, and runtime
monitoring. Communication latency and bandwidth limitations
can significantly influence migration overhead and overall
system schedulability. The co-design of communication
scheduling policies and task scheduling policies for
distributed real-time systems with migration support
constitutes one of the open research problems motivating
the work presented in this dissertation.

% ============================================================
\section{Summary}
\label{sec:background_summary}

This chapter introduced the fundamental concepts relevant to
the research presented in this dissertation. Starting from
the foundations of real-time systems, scheduling theory,
and operating systems support, the discussion progressed to
mixed-criticality models and their associated challenges in
resource sharing, mode management, and safety certification.
The chapter then addressed distributed real-time system
architectures, covering end-to-end timing analysis, the
automotive evolution towards centralized and zonal platforms,
heterogeneous and edge computing paradigms, and the specific
challenges of task migration in real-time contexts. Finally,
relevant communication technologies---from legacy fieldbuses
to TSN and service-oriented middleware---were reviewed.

The presented concepts collectively establish the theoretical
foundation required for understanding the state of the art
and the research challenges discussed in the following
chapter. In particular, the intersecting requirements of
mixed-criticality scheduling, distributed resource management,
and deterministic communication under task migration scenarios
motivate the design space explored in the remainder of this
dissertation.

\chapter{Related Work and State of the Art}
\label{chap:relatedwork}

The increasing complexity of modern embedded and cyber-physical
systems has led to extensive research on architectures that
combine deterministic timing guarantees with flexibility and
runtime adaptability. Classical real-time systems were designed
under static assumptions regarding task sets, execution platforms,
and scheduling policies. However, modern automotive and
distributed embedded systems increasingly require dynamic
adaptation to workload variations, heterogeneous hardware, and
runtime reconfiguration.

This chapter reviews the most relevant areas of research related
to this dissertation, including real-time scheduling, mixed-
criticality systems, distributed embedded architectures, task
migration, and adaptive resource management.

% ============================================================
\section{Scheduling and Resource Management in Real-Time Systems}

Real-time systems require that computational tasks not only
produce correct results, but also meet strict timing constraints.
Foundational work in real-time scheduling established classical
models for periodic task systems and provided schedulability
conditions for fixed-priority and dynamic-priority scheduling
algorithms~\cite{liu1973,buttazzo2011}.

Rate Monotonic (RM) and Earliest Deadline First (EDF) scheduling
remain fundamental paradigms for real-time task scheduling.
However, their applicability becomes limited in multicore and
distributed environments due to migration costs, shared resource
contention, and timing unpredictability.

\subsection{Multiprocessor and Distributed Scheduling}

Multiprocessor real-time scheduling introduces additional
challenges such as task migration overhead, cache effects,
and synchronization costs. Megel et al.\ investigated
migration-aware scheduling strategies aimed at minimizing
preemption and migration overhead in multiprocessor systems
~\cite{megel2010}. However, their work remains primarily
theoretical and does not address runtime distributed execution
across heterogeneous nodes.

Recent research trends highlight that scheduling decisions
can no longer be considered independently from communication
systems, particularly in distributed embedded systems where
network latency is part of end-to-end timing guarantees.

% ============================================================
\section{Mixed-Criticality Systems}

Mixed-criticality systems (MCS) enable the execution of tasks
with different safety and timing requirements on shared platforms.
This paradigm is widely used in automotive and avionics domains
where safety-critical and non-critical workloads coexist.

The foundational model by Vestal introduced multiple WCET
assumptions depending on certification levels, enabling more
flexible scheduling strategies~\cite{vestal2007}. Since then,
extensive research has focused on improving schedulability
analysis and runtime behavior.

Burns and Davis provide a comprehensive survey of mixed-criticality
systems and highlight the trade-off between resource utilization
and timing guarantees~\cite{burns2017}. Traditional approaches
typically rely on mode-switching mechanisms, where low-criticality
tasks are degraded or dropped under overload conditions.

Recent work explores more adaptive strategies, including energy-
aware and multicore-aware scheduling approaches~\cite{mcsEnergySurvey2022}.
However, most approaches still assume static task allocation and
do not support runtime migration across distributed systems.

% ============================================================
\section{Distributed Real-Time Systems and Automotive Architectures}

Distributed real-time systems consist of multiple computational
nodes connected via a communication network. These systems are
widely used in automotive, aerospace, and industrial automation.

Modern automotive architectures are transitioning from distributed
Electronic Control Unit (ECU)-based systems towards centralized
and zonal architectures. This transition is driven by increasing
computational demands from advanced driver assistance systems
and autonomous driving functionalities~\cite{baunach2018}.

Recent developments in software-defined vehicles (SDVs) and
centralized automotive computing platforms further accelerate
this transformation~\cite{sdvSurvey2024}. These architectures
aim to consolidate computation onto high-performance platforms
while maintaining real-time guarantees.

At the same time, Time-Sensitive Networking (TSN) has emerged
as a key enabling technology for deterministic communication in
Ethernet-based systems. TSN introduces mechanisms such as the
Time-Aware Shaper (TAS), which enables time-triggered scheduling
of network traffic.

A comprehensive survey of TSN scheduling mechanisms is provided
by Stüber et al.~\cite{tsnSurvey2023}, while more recent work
investigates experimental evaluations of TSN scheduling approaches
in industrial scenarios~\cite{tsnRtas2024}.

Despite these advances, most distributed real-time architectures
still rely on static task allocation, where execution placement
is determined offline and remains unchanged during runtime.

\section{Task Offloading and Orchestration in IoT and Edge Systems}

In parallel to research on distributed real-time systems, a large
body of work in the Internet of Things (IoT), edge computing, and
fog computing has investigated the problem of task offloading and
resource orchestration across heterogeneous computing nodes.

In IoT and edge environments, resource-constrained devices often
lack sufficient computational capacity to execute complex workloads
locally. To address this limitation, task offloading mechanisms
transfer computation to more capable edge or cloud resources,
reducing execution latency, energy consumption, or improving
system scalability.

Mao et al.\ studied computation offloading strategies in mobile
edge computing systems and highlighted the trade-off between
communication cost, computation latency, and energy consumption
\cite{mao2017mobile}. Their work shows that optimal offloading
decisions depend on dynamic system conditions such as network
load and device state.

Similarly, You et al.\ analyzed joint communication and computation
resource allocation for mobile edge systems, demonstrating that
offloading decisions must consider both network and compute
constraints simultaneously \cite{you2017energy}.

Recent research in fog and edge computing has extended these ideas
towards multi-node orchestration frameworks, where workloads are
dynamically distributed across hierarchical computing layers
(edge, fog, cloud). These systems often rely on centralized or
semi-centralized orchestration components that continuously
optimize task placement based on system state information.

Shi et al.\ provide a comprehensive survey of edge computing
architectures and highlight orchestration as a key challenge for
scalable IoT systems \cite{shi2016edge}.

More recent studies emphasize the use of container orchestration
frameworks and serverless computing models to dynamically manage
task placement in distributed environments. However, these systems
are typically designed for best-effort execution environments and
do not provide hard real-time guarantees.

While IoT and edge computing research provides powerful models for
dynamic task offloading and system-wide orchestration, most existing
approaches assume relaxed timing constraints and prioritize
throughput, energy efficiency, or cost optimization.

In contrast, real-time embedded systems require deterministic timing
guarantees, bounded execution latency, and predictable scheduling
behavior. This fundamental difference limits the direct applicability
of IoT offloading techniques to safety-critical real-time systems.

Nevertheless, the conceptual overlap between task offloading in IoT
systems and task migration in distributed real-time systems provides
an important foundation for this dissertation. Both domains address
the problem of dynamically assigning computational tasks to
distributed resources under changing system conditions, albeit with
different constraints and objectives.

% ============================================================
\section{Task Migration in Real-Time Systems}

Task migration has been widely studied as a mechanism to improve
load balancing, fault tolerance, and resource utilization in
real-time systems. However, it introduces significant challenges
related to timing predictability, communication overhead, and
state consistency.

Migration-aware scheduling approaches primarily aim at reducing
migration overhead while preserving schedulability guarantees.
Megel et al.\ demonstrated the cost sensitivity of migration in
multiprocessor scheduling~\cite{megel2010}.

In fault-tolerant systems, migration is often used for recovery
purposes. Chen and Lyu proposed a FlexRay-based framework in which
backup ECUs take over execution in case of failure~\cite{chen2012}.
However, such approaches operate at coarse granularity (ECU level)
and do not support fine-grained task migration between active nodes.

Recent research in TSN and distributed systems highlights that
dynamic routing and scheduling are tightly coupled problems,
further increasing the complexity of runtime adaptation
~\cite{tsnSchedAnalysis2024}.

Overall, most existing approaches either focus on theoretical
scheduling models or coarse-grained fault recovery mechanisms,
with limited support for runtime task migration in heterogeneous
distributed systems.

% ============================================================
\section{Machine Learning for Scheduling and Deployment}

Machine learning techniques have been increasingly applied to
scheduling and resource allocation problems due to their ability
to approximate complex optimization problems.

Early work explored neural network approaches for schedulability
analysis and real-time scheduling prediction~\cite{cardeira1994,dominguez2004}.
These approaches primarily focused on offline system design.

More recent research investigates learning-based approaches for
resource allocation in dynamic systems. In particular, reinforcement
learning has been applied to distributed scheduling and routing
problems, especially in TSN and edge environments.

However, despite promising results, most machine learning-based
approaches remain limited to simulation or offline optimization
and are not integrated into real-time operating systems with
strict deterministic guarantees.

% ============================================================
\section{Existing Architectures and Frameworks}

Several research projects have investigated architectures for
adaptive deployment and distributed execution in automotive systems.

The KIA4SM project proposed a cooperative integration architecture
for future smart mobility systems, enabling task migration across
heterogeneous automotive platforms and introducing checkpoint-restore
mechanisms~\cite{eckl2015}.

While such architectures demonstrate the feasibility of distributed
execution, they remain primarily conceptual or prototype-level
systems without full integration into real-time operating systems
supporting fine-grained runtime migration.

Recent research in virtualization and partitioning further explores
isolation mechanisms for mixed-criticality systems, but these
approaches typically assume static deployment models and do not
support dynamic task relocation.

% ============================================================
\section{Research Gap}

The reviewed literature shows significant progress in real-time
scheduling, mixed-criticality systems, distributed architectures,
and task migration. However, important limitations remain.

Most existing approaches treat scheduling, communication, and
deployment as separate concerns. Real-time systems typically rely
on static task allocation, while mixed-criticality systems focus on
controlled degradation rather than runtime adaptation. Distributed
architectures improve scalability but lack fine-grained dynamic
reconfiguration capabilities.

Task migration is generally addressed either in theoretical models
or coarse-grained fault-tolerant systems. Machine learning approaches
provide promising optimization capabilities but are rarely integrated
into safety-critical real-time execution environments.

Consequently, there is a lack of integrated system-level architectures
capable of supporting fine-grained runtime task migration across
distributed heterogeneous mixed-criticality systems while preserving
strict timing guarantees.

This dissertation addresses this gap by proposing a distributed
real-time operating system architecture enabling adaptive task
migration across heterogeneous execution platforms.

%\begin{tikzpicture}[scale=1.1]
%	
%	% Axes
%	\draw[->, thick] (0,0) -- (10,0) node[right] {Adaptability};
%	\draw[->, thick] (0,0) -- (0,8) node[above] {System Scope};
%	
%	% Axis labels
%	\node at (2,-0.7) {Static};
%	\node at (8,-0.7) {Dynamic};
%	
%	\node[rotate=90] at (-0.8,2) {Single Node};
%	\node[rotate=90] at (-0.8,6.5) {Distributed};
%	
%	% Regions / Categories
%	
%	% Static RT Systems
%	\node[draw, fill=gray!20, rounded corners] at (2,1.5) {
%		\small Static RT Systems
%	};
%	
%	% Mixed-Criticality Systems
%	\node[draw, fill=blue!20, rounded corners] at (3.5,2.5) {
%		\small Mixed-Criticality
%	};
%	
%	% Distributed Static Systems
%	\node[draw, fill=green!20, rounded corners] at (3,5.5) {
%		\small Distributed RT\\(Static Mapping)
%	};
%	
%	% Migration / Adaptive Systems
%	\node[draw, fill=orange!20, rounded corners] at (7,5) {
%		\small Migration-Based\\Approaches
%	};
%	
%	% Emerging Adaptive Systems
%	\node[draw, fill=yellow!30, rounded corners] at (7.5,6.5) {
%		\small Adaptive Distributed\\Systems
%	};
%	
%	% Your Contribution
%	\node[draw, fill=red!30, thick, rounded corners] at (8.5,7.2) {
%		\small \textbf{Your Approach}\\
%		\small (Dynamic MC Migration)
%	};
%	
%	% Optional arrows (evolution)
%	\draw[->, dashed] (2,1.5) -- (3.5,2.5);
%	\draw[->, dashed] (3.5,2.5) -- (3,5.5);
%	\draw[->, dashed] (3,5.5) -- (7,5);
%	\draw[->, dashed] (7,5) -- (8.5,7.2);
%	
%\end{tikzpicture}
